\chapter{2018年全国II卷（理）}

\section{单选题}

\begin{problem}
\(\dfrac{1+2\i}{1-2\i}=\)（\quad）

\choices
    {\(-\dfrac35-\dfrac45\i\)}
    {\(\dfrac35-\dfrac45\i\)}
    {\(-\dfrac35+\dfrac45\i\)}
    {\(\dfrac35+\dfrac45\i\)}

\end{problem}

\begin{answer}
C
\end{answer}

\begin{solution}
利用复数除法，\(\frac{1+2\i}{1-2\i}=\frac{(1+2\i)(1+2\i)}{(1-2\i)(1+2\i)}=\frac{1+4\i+4\i^2}{1+4}=-\frac35+\frac45\i\). 故选 C.
\end{solution}

\begin{problem}
已知集合 \(A=\{(x,y)\mid x^2+y^2\le 3,\ x\in\Z,\ y\in\Z\}\)，则 \(A\) 中元素的个数为（\quad）

\choices
    {\(9\)}
    {\(8\)}
    {\(5\)}
    {\(4\)}

\end{problem}

\begin{answer}
A
\end{answer}

\begin{solution}
因为 \(x\)，\(y\) 都是整数且 \(x^2+y^2\le 3\)，所以 \(x\) 只能取 \(-1\)，\(0\)，\(1\). 分别讨论可得：
\(x=-1\text{ 时， }y=-1\)，\(0\)，\(1;\)
\(x=0\text{ 时， }y=-1\)，\(0\)，\(1;\)
\(x=1\text{ 时， }y=-1\)，\(0\)，\(1\).
因此 A 中元素共有 \(9\) 个，故选 A.
\end{solution}

\begin{problem}
函数 \(f(x)=\dfrac{e^x-e^{-x}}{x^2}\) 的图象大致为（\quad）

\choices
    {\choicebitmap[width=0.15\paperwidth]{img/2018/national_paper_2_science/q3_fig1.png}}
    {\choicebitmap[width=0.15\paperwidth]{img/2018/national_paper_2_science/q3_fig2.png}}
    {\choicebitmap[width=0.15\paperwidth]{img/2018/national_paper_2_science/q3_fig3.png}}
    {\choicebitmap[width=0.15\paperwidth]{img/2018/national_paper_2_science/q3_fig4.png}}

\end{problem}

\begin{answer}
B
\end{answer}

\begin{solution}
函数 \(f(x)\) 是奇函数，图象关于原点对称；且当 \(x\to +\infty\) 时，
\(f(x)=\frac{e^x-e^{-x}}{x^2}\to +\infty\)，
同时在 \(x>0\) 时函数为正. 结合选项图象可知选 B.
\end{solution}

\begin{problem}
已知向量 \(\bs a\)，\(\bs b\) 满足 \(|\bs a|=1\)，\(\bs a\cdot\bs b=-1\)，则 \(\bs a\cdot(2\bs a-\bs b)=\)（\quad）

\choices
    {\(4\)}
    {\(3\)}
    {\(2\)}
    {\(0\)}

\end{problem}

\begin{answer}
B
\end{answer}

\begin{solution}
\(\bs a\cdot(2\bs a-\bs b)=2|\bs a|^2-\bs a\cdot\bs b=2-(-1)=3\).
故选 B.
\end{solution}

\begin{problem}
双曲线 \(\dfrac{x^2}{a^2}-\dfrac{y^2}{b^2}=1\)（\(a>0\)，\(b>0\)）的离心率为 \(\sqrt3\)，则其渐近线方程为（\quad）

\choices
    {\(y=\pm\sqrt2\,x\)}
    {\(y=\pm\sqrt3\,x\)}
    {\(y=\pm\dfrac{\sqrt2}{2}x\)}
    {\(y=\pm\dfrac{\sqrt3}{2}x\)}

\end{problem}

\begin{answer}
A
\end{answer}

\begin{solution}
由双曲线离心率 \(e=\dfrac ca=\sqrt3\)，得
\(\frac{b^2}{a^2}=e^2-1=2\)，
所以渐近线方程为
\(y=\pm\frac ba x=\pm\sqrt2\,x\).
故选 A.
\end{solution}

\begin{problem}
在 \(\triangle ABC\) 中，\(\cos A=\dfrac45\)，\(BC=3\)，\(AC=5\)，则 \(AB=\)（\quad）

\choices
    {\(4\)}
    {\(\sqrt{19}\)}
    {\(3\)}
    {\(2\)}

\end{problem}

\begin{answer}
A
\end{answer}

\begin{solution}
由余弦定理，
\(BC^2=AB^2+AC^2-2\cdot AB\cdot AC\cos A\).
代入 \(BC=3\)，\(AC=5\)，\(\cos A=\dfrac45\)，得
\(9=AB^2+25-8AB\).
整理得
\(AB^2-8AB+16=0\)，
所以
\(AB=4\).
故选 A.
\end{solution}

\begin{problem}
为计算 \(S=1-\dfrac12+\dfrac13-\dfrac14+\cdots+\dfrac1{100}\)，设计了如图的程序框图，则在空白框中应填入（\quad）

\texfigure[max width=0.38\linewidth]{img_repaint/2018/national_paper_2_science/q7_fig1.tex}

\choices
    {\(i=i+1\)}
    {\(i=i+2\)}
    {\(i=i+3\)}
    {\(i=i+4\)}

\end{problem}

\begin{answer}
B
\end{answer}

\begin{solution}
程序每次累加后都跳过一个数，因此步长为 \(2\)，空白框中应填入
\(i=i+2\).
故选 B.
\end{solution}

\begin{problem}
我国数学家陈景润在哥德巴赫猜想的研究中取得了世界领先的成果. 哥德巴赫猜想是“每个大于 \(2\) 的偶数可以表示为两个素数的和”，如 \(30=7+23\). 在不超过 \(30\) 的素数中，随机选取两个不同的数，其和等于 \(30\) 的概率是（\quad）

\choices
    {\(\dfrac1{12}\)}
    {\(\dfrac1{14}\)}
    {\(\dfrac1{15}\)}
    {\(\dfrac1{18}\)}

\end{problem}

\begin{answer}
C
\end{answer}

\begin{solution}
不超过 \(30\) 的素数有
\(2\)，\(3\)，\(5\)，\(7\)，\(11\)，\(13\)，\(17\)，\(19\)，\(23\)，\(29\)，
共 \(10\) 个. 从中选 \(2\) 个不同的数共有 \(\binom{10}{2}=45\) 种，其中和为 \(30\) 的有
\((7,23)\)，\((11,19)\)，\((13,17)\)，
共 \(3\) 种，所以所求概率为
\(\frac3{45}=\frac1{15}\).
故选 C.
\end{solution}

\begin{problem}
在长方体 \(ABCD-A_1B_1C_1D_1\) 中，\(AB=BC=1\)，\(AA_1=\sqrt3\)，则异面直线 \(AD_1\) 与 \(DB_1\) 所成角的余弦值为（\quad）.

\choices
    {\(\dfrac15\)}
    {\(\dfrac{\sqrt5}{6}\)}
    {\(\dfrac{\sqrt5}{5}\)}
    {\(\dfrac{\sqrt2}{2}\)}

\end{problem}

\begin{answer}
C
\end{answer}

\begin{solution}
设 \(A(0,0,0)\)，\(B(1,0,0)\)，\(D(0,1,0)\)，\(A_1(0,0,\sqrt3)\)，则
\(D_1(0,1,\sqrt3)\)，\(B_1(1,0,\sqrt3)\).
于是
\(\overrightarrow{AD_1}=(0,1,\sqrt3)\)，\(\overrightarrow{DB_1}=(1,-1,\sqrt3)\).
故
\(\overrightarrow{AD_1}\cdot\overrightarrow{DB_1}=0-1+3=2\)，
\(\lVert \overrightarrow{AD_1}\rVert=2\)，\(\lVert \overrightarrow{DB_1}\rVert=\sqrt5\).
所以所成角的余弦值为
\(\frac{2}{2\sqrt5}=\frac{\sqrt5}{5}\).
故选 C.
\end{solution}

\begin{problem}
若 \(f(x)=\cos x-\sin x\) 在 \([-a,a]\) 是减函数，则 \(a\) 的最大值是（\quad）

\choices
    {\(\dfrac{\pi}{4}\)}
    {\(\dfrac{\pi}{2}\)}
    {\(\dfrac{3\pi}{4}\)}
    {\(\pi\)}

\end{problem}

\begin{answer}
A
\end{answer}

\begin{solution}
\(f(x)=\cos x-\sin x=\sqrt2\cos\!\left(x+\frac{\pi}{4}\right)\).
当
\(x+\frac{\pi}{4}\in[0,\pi]\)
时，\(f(x)\) 单调递减，即
\(x\in\left[-\frac{\pi}{4},\frac{3\pi}{4}\right]\).
要使 \([-a,a]\) 仍处于该减区间内，最大的 \(a\) 为
\(a=\frac{\pi}{4}\).
故选 A.
\end{solution}

\begin{problem}
已知 \(f(x)\) 是定义域为 \((-\infty,+\infty)\) 的奇函数，满足 \(f(1-x)=f(1+x)\)，若 \(f(1)=2\)，则 \(f(1)+f(2)+f(3)+\cdots+f(50)=\)（\quad）

\choices
    {\(-50\)}
    {\(0\)}
    {\(2\)}
    {\(50\)}

\end{problem}

\begin{answer}
C
\end{answer}

\begin{solution}
由 \(f(1-x)=f(1+x)\)，令 \(x\mapsto x-1\)，得
\(f(-x)=f(x+2)\).
因为 \(f(x)\) 是奇函数，所以
\(f(x+2)=-f(x)\)，
从而
\(f(x+4)=f(x)\).
于是
\(f(1)=2\)，\(f(2)=0\)，\(f(3)=-2\)，\(f(4)=0\)，
之后按 \(4\) 个一循环. 故
\(\sum_{k=1}^{50}f(k)=12\cdot(2+0-2+0)+f(49)+f(50)=2\).
故选 C.
\end{solution}

\begin{problem}
已知 \(F_1\)，\(F_2\) 是椭圆 \(C:\dfrac{x^2}{a^2}+\dfrac{y^2}{b^2}=1\)（\(a>b>0\)）的左，右焦点，\(A\) 是 \(C\) 的左顶点，点 \(P\) 在过 \(A\) 且斜率为 \(\dfrac{\sqrt3}{6}\) 的直线上，\(\triangle PF_1F_2\) 为等腰三角形，\(\angle F_1F_2P=120^\circ\)，则 \(C\) 的离心率为（\quad）

\choices
    {\(\dfrac23\)}
    {\(\dfrac12\)}
    {\(\dfrac13\)}
    {\(\dfrac14\)}

\end{problem}

\begin{answer}
D
\end{answer}

\begin{solution}
设椭圆焦距参数为 \(c\)，则
\(F_1(-c,0)\)，\(F_2(c,0)\)，\(A(-a,0)\).
由 \(\angle F_1F_2P=120^\circ\) 且 \(\triangle PF_1F_2\) 为等腰三角形，可得
\(PF_2=F_1F_2=2c\).
于是点 \(P\) 的坐标为
\((2c,\sqrt3\,c)\).
又 \(P\) 在过 \(A\) 且斜率为 \(\dfrac{\sqrt3}{6}\) 的直线上，所以
\(\frac{\sqrt3\,c}{2c+a}=\frac{\sqrt3}{6}\).
解得
\(a=4c\).
因此椭圆离心率
\(e=\frac ca=\frac14\).
故选 D.
\end{solution}
\section{填空题}

\begin{problem}
曲线 \(y=2\ln(x+1)\) 在点 \((0,0)\) 处的切线方程为\fillinblank{}.

\end{problem}

\begin{answer}
\(y=2x\)
\end{answer}

\begin{solution}
设 \(g(x)=2\ln(x+1)\)，则
\(g'(x)=\frac2{x+1}\)，
所以切线斜率为 \(g'(0)=2\). 又切点为 \((0,0)\)，故切线方程为
\(y=2x\).
\end{solution}

\begin{problem}
若 \(x\)，\(y\) 满足约束条件 \((2y-x-3\le 0,\ x+2y-5\ge 0,\ x\le 5)\)，则 \(z=x+y\) 的最大值为\fillinblank{}.

\bitmapfigure[max width=0.42\linewidth]{img/2018/national_paper_2_science/q14_fig1.png}

\end{problem}

\begin{answer}
\(9\)
\end{answer}

\begin{solution}
作出可行域如图所示，目标函数 \(z=x+y\) 的等值线为 \(y=-x+z\). 平移该直线，最先经过可行域顶点 \(A(5,4)\) 时取到最大值，
\(z_{\max}=5+4=9\).
故答案为 \(9\).
\end{solution}

\begin{problem}
已知 \(\sin\alpha+\cos\beta=1\)，\(\cos\alpha+\sin\beta=0\)，则 \(\sin(\alpha+\beta)=\fillinblank{}\).

\end{problem}

\begin{answer}
\(-\dfrac12\)
\end{answer}

\begin{solution}
由
\(\sin\alpha+\cos\beta=1\)，
两边平方得
\(\sin^2\alpha+2\sin\alpha\cos\beta+\cos^2\beta=1\).
由
\(\cos\alpha+\sin\beta=0\)，
两边平方得
\(\cos^2\alpha+2\cos\alpha\sin\beta+\sin^2\beta=0\).
将上述两式相加，得
\(2+2(\sin\alpha\cos\beta+\cos\alpha\sin\beta)=1\)，
即
\(2\sin(\alpha+\beta)=-1\).
所以
\(\sin(\alpha+\beta)=-\frac12\).
\end{solution}

\begin{problem}
已知圆锥的顶点为 \(S\)，母线 \(SA\)，\(SB\) 所成角的余弦值为 \(\dfrac78\)，\(SA\) 与圆锥底面所成角为 \(45^\circ\)，若 \(\triangle SAB\) 的面积为 \(5\sqrt{15}\)，则该圆锥的侧面积为\fillinblank{}.

\end{problem}

\begin{answer}
\(40\sqrt2\pi\)
\end{answer}

\begin{solution}
设圆锥母线长为 \(l\). 由
\(\cos\angle ASB=\frac78\)
得
\(\sin\angle ASB=\frac{\sqrt{15}}8\).
又
\(S_{\triangle SAB}=\frac12 l^2\sin\angle ASB=5\sqrt{15}\)，
所以
\(\frac12 l^2\cdot\frac{\sqrt{15}}8=5\sqrt{15}\)，
解得
\(l^2=80\)，\(l=4\sqrt5\).
因为 \(SA\) 与底面所成角为 \(45^\circ\)，所以底面半径
\(r=l\sin45^\circ=2\sqrt{10}\).
故圆锥侧面积为
\(S=\pi rl=\pi\cdot 2\sqrt{10}\cdot 4\sqrt5=40\sqrt2\pi\).
故答案为 \(40\sqrt2\pi\).
\end{solution}
\section{解答题}

\begin{problem}
记 \(S_n\) 为等差数列 \(\{a_n\}\) 的前 \(n\) 项和，已知 \(a_1=-7\)，\(S_3=-15\).

\begin{enumerate}
\item 求 \(\{a_n\}\) 的通项公式；
\item 求 \(S_n\)，并求 \(S_n\) 的最小值.
\end{enumerate}

\end{problem}

\begin{solution}
（1）因为等差数列 \(\{a_n\}\) 中，\(a_1=-7\)，\(S_3=-15\)，
\(a_1=-7\)，\(3a_1+3d=-15\)，
解得
\(a_1=-7\)，\(d=2\).
所以
\(a_n=-7+2(n-1)=2n-9\).

（2）因为 \(a_1=-7\)，\(d=2\)，\(a_n=2n-9\)，
\(S_n=\frac n2(a_1+a_n)=\frac12\bigl(2n^2-16n\bigr)=n^2-8n=(n-4)^2-16\).
所以当 \(n=4\) 时，前 \(n\) 项和 \(S_n\) 取得最小值 \(-16\).
\end{solution}

\begin{problem}
如图是某地区 \(2000\) 年至 \(2016\) 年环境基础设施投资额 \(y\)（单位：亿元）的折线图.

为了预测该地区 \(2018\) 年的环境基础设施投资额，建立了 \(y\) 与时间变量 \(t\) 的两个线性回归模型. 根据 \(2000\) 年至 \(2016\) 年的数据（时间变量 \(t\) 的值依次为 \(1\)，\(2\)，\(\cdots\)，\(17\)）建立模型（1）：\(\hat y=-30.4+13.5t\)；根据 \(2010\) 年至 \(2016\) 年的数据（时间变量 \(t\) 的值依次为 \(1\)，\(2\)，\(\cdots\)，\(7\)）建立模型（2）：\(\hat y=99+17.5t\).

\begin{enumerate}
\item 分别利用这两个模型，求该地区 \(2018\) 年的环境基础设施投资额的预测值；
\item 你认为用哪个模型得到的预测值更可靠？并说明理由.
\end{enumerate}

\bitmapfigure[max width=0.84\linewidth]{img/2018/national_paper_2_science/q18_fig1.png}

\end{problem}

\begin{solution}
（1）根据模型（1）
\(\hat y=-30.4+13.5t\)，
计算 \(t=19\) 时，
\(\hat y=-30.4+13.5\times 19=226.1\).
利用这个模型，求出该地区 \(2018\) 年的环境基础设施投资额的预测值是 \(226.1\) 亿元.
根据模型（2）
\(\hat y=99+17.5t\)，
计算 \(t=9\) 时，
\(\hat y=99+17.5\times 9=256.5\).
利用这个模型，求出该地区 \(2018\) 年的环境基础设施投资额的预测值是 \(256.5\) 亿元.

（2）模型（2）得到的预测值更可靠.

因为从总体数据看，该地区从 \(2000\) 年到 \(2016\) 年的环境基础设施投资额是逐年上升的，而从 \(2000\) 年到 \(2009\) 年间递增的幅度较小些，从 \(2010\) 年到 \(2016\) 年间递增的幅度较大些，所以利用模型（2）的预测值更可靠些.
\end{solution}

\begin{problem}
设抛物线 \(C:y^2=4x\) 的焦点为 \(F\)，过 \(F\) 且斜率为 \(k\)（\(k>0\)）的直线 \(l\) 与 \(C\) 交于 \(A\)，\(B\) 两点，\(|AB|=8\).

\begin{enumerate}
\item 求 \(l\) 的方程；
\item 求过点 \(A\)，\(B\) 且与 \(C\) 的准线相切的圆的方程.
\end{enumerate}

\end{problem}

\begin{solution}
（1）【法一】抛物线 \(C:y^2=4x\) 的焦点为 \(F(1,0)\).
设直线 \(AB\) 的方程为
\(y=k(x-1)\)，
设 \(A(x_1,y_1)\)，\(B(x_2,y_2)\)，则
\examdisplaycases{}{
y=k(x-1),\\
y^2=4x
}
整理得
\(k^2x^2-2(k^2+2)x+k^2=0\)，
所以
\(x_1+x_2=\frac{2(k^2+2)}{k^2}\)，\(x_1x_2=1\).
由抛物线的焦点弦公式，
\(|AB|=x_1+x_2+p=\frac{2(k^2+2)}{k^2}+2\).
因为 \(|AB|=8\)，解得
\(k^2=1\).
又 \(k>0\)，所以
\(k=1\).
因此直线 \(l\) 的方程为
\(y=x-1\).

【法二】抛物线 \(C:y^2=4x\) 的焦点为 \(F(1,0)\)，设直线 \(AB\) 的倾斜角为 \(\theta\). 由抛物线的焦点弦公式，
\(|AB|=\frac{2p}{\sin^2\theta}=\frac{4}{\sin^2\theta}=8\)，
解得
\(\sin^2\theta=\frac12\).
因为 \(k>0\)，所以
\(\theta=\frac{\pi}{4}\)，
从而
\(k=1\).
因此直线 \(l\) 的方程为
\(y=x-1\).

（2）由（1）可得 \(AB\) 的中点坐标为 \(D(3,2)\)，则直线 \(AB\) 的垂直平分线方程为
\(y-2=-(x-3)\)，
即
\(y=-x+5\).
设所求圆的圆心坐标为 \((x_0,y_0)\)，则
\(y_0=-x_0+5\).
又圆与准线 \(x=-1\) 相切，所以半径为
\(r=x_0+1\).
由抛物线的定义和圆经过 \(A\)，\(B\) 两点，可得
\((x_0+1)^2=\frac{(y_0-x_0+1)^2}{2}+16\).
联立
\(y_0=-x_0+5\)
与上式，解得
\(x_0=3\)，\(y_0=2\)
或
\(x_0=11\)，\(y_0=-6\).
因此，所求圆的方程为
\((x-3)^2+(y-2)^2=16\)
或
\((x-11)^2+(y+6)^2=144\).

\solutionfigure{\bitmapfigure[max width=0.80\linewidth]{img/2018/national_paper_2_science/q19_solution_fig1.png}}
\end{solution}

\begin{problem}
如图，在三棱锥 \(P-ABC\) 中，\(AB=BC=2\sqrt2\)，\(PA=PB=PC=AC=4\)，\(O\) 为 \(AC\) 的中点.

\begin{enumerate}
\item 证明：\(PO\perp\) 平面 \(ABC\)；
\item 若点 \(M\) 在棱 \(BC\) 上，且二面角 \(M-PA-C\) 为 \(30^\circ\)，求 \(PC\) 与平面 \(PAM\) 所成角的正弦值.
\end{enumerate}

\bitmapfigure[max width=0.38\linewidth]{img/2018/national_paper_2_science/q20_fig1.png}

\end{problem}

\begin{solution}
（1）证明：连接 \(BO\).

因为 \(AB=BC=2\sqrt2\)，\(O\) 是 \(AC\) 的中点，所以
\(BO\perp AC\)，\(BO=2\).
又 \(PA=PB=PC=AC=4\)，所以
\(PO\perp AC\)，\(PO=2\sqrt3\).
由
\(PB^2=PO^2+BO^2\)
可得
\(PO\perp OB\).
因为
\(OB\cap AC=O\)，
所以
\(PO\perp \text{平面 }ABC\).

（2）建立以 \(O\) 为原点，\(OB\)，\(OC\)，\(OP\) 分别为 \(x\)，\(y\)，\(z\) 轴的空间直角坐标系，则
\(A(0,-2,0)\)，\(P(0,0,2\sqrt3)\)，\(C(0,2,0)\)，\(B(2,0,0)\).
于是
\(\overrightarrow{BC}=(-2,2,0)\).
设
\(\overrightarrow{BM}=\lambda \overrightarrow{BC}=(-2\lambda,2\lambda,0)\)，\(0<\lambda<1\).
则
\(\overrightarrow{AM}=\overrightarrow{BM}-\overrightarrow{BA}=(-2\lambda,2\lambda,0)-(-2,-2,0)=(2-2\lambda,2\lambda+2,0)\).
平面 \(PAC\) 的法向量为
\(\bs m=(1,0,0)\).
设平面 \(MPA\) 的法向量为
\(\bs n=(x,y,z)\)，
则
\(\overrightarrow{PA}=(0,-2,-2\sqrt3)\).
于是
\(\bs n\cdot \overrightarrow{PA}=-2y-2\sqrt3\,z=0\)，
\(\bs n\cdot \overrightarrow{AM}=(2-2\lambda)x+(2\lambda+2)y=0\).
令 \(z=1\)，则
\(y=-\sqrt3\)，\(x=\frac{(\lambda+1)\sqrt3}{1-\lambda}\).
所以可取
\(\bs n=\left(\frac{(\lambda+1)\sqrt3}{1-\lambda},-\sqrt3,1\right)\).
因为二面角 \(M-PA-C\) 为 \(30^\circ\)，所以
\(\cos30^\circ=\frac{|\bs m\cdot \bs n|}{|\bs m||\bs n|}=\frac{\sqrt3}{2}\).
即
\(\frac{\left|\dfrac{(\lambda+1)\sqrt3}{1-\lambda}\right|}{\sqrt{\left(\dfrac{(\lambda+1)\sqrt3}{1-\lambda}\right)^2+(-\sqrt3)^2+1}}=\frac{\sqrt3}{2}\).
解得
\(\lambda=\frac13\)
或
\(\lambda=3\).
因为 \(0<\lambda<1\)，所以
\(\lambda=\frac13\).
于是平面 \(MPA\) 的法向量可取为
\(\bs n=(2\sqrt3,-\sqrt3,1)\).

又
\(\overrightarrow{PC}=(0,2,-2\sqrt3)\)，
所以 \(PC\) 与平面 \(PAM\) 所成角的正弦值为
\(\sin\theta=\left|\cos\langle \overrightarrow{PC},\bs n\rangle\right|=\frac{|\overrightarrow{PC}\cdot \bs n|}{|\overrightarrow{PC}||\bs n|}=\frac{\sqrt3}{4}\).

\solutionfigure{\bitmapfigure[max width=0.38\linewidth]{img/2018/national_paper_2_science/q20_solution_fig1.png}}
\end{solution}

\begin{problem}
已知函数 \(f(x)=e^x-ax^2\).

\begin{enumerate}
\item 若 \(a=1\)，证明：当 \(x\ge 0\) 时，\(f(x)\ge 1\)；
\item 若 \(f(x)\) 在 \((0,+\infty)\) 只有一个零点，求 \(a\).
\end{enumerate}

\end{problem}

\begin{solution}
（1）当 \(a=1\) 时，函数
\(f(x)=e^x-x^2\).
则
\(f'(x)=e^x-2x\).
令
\(g(x)=e^x-2x\)，
则
\(g'(x)=e^x-2\).
令 \(g'(x)=0\)，得
\(x=\ln2\).
当 \(x\in(0,\ln2)\) 时，\(g'(x)<0\)；当 \(x\in(\ln2,+\infty)\) 时，\(g'(x)>0\).
所以
\(g(x)\ge g(\ln2)=e^{\ln2}-2\ln2=2-2\ln2>0\).
因此 \(f(x)\) 在 \([0,+\infty)\) 单调递增，从而
\(f(x)\ge f(0)=1\).

（2）【法一】\(f(x)\) 在 \((0,+\infty)\) 只有一个零点，等价于方程
\(e^x-ax^2=0\)
在 \((0,+\infty)\) 只有一个根，也等价于
\(a=\frac{e^x}{x^2}\)
在 \((0,+\infty)\) 只有一个根，即函数 \(y=a\) 与
\(G(x)=\frac{e^x}{x^2}\)
的图象在 \((0,+\infty)\) 只有一个交点.

因为
\(G'(x)=\frac{e^x(x-2)}{x^3}\)，
所以当 \(x\in(0,2)\) 时，\(G'(x)<0\)；当 \(x\in(2,+\infty)\) 时，\(G'(x)>0\).
因此 \(G(x)\) 在 \((0,2)\) 递减，在 \((2,+\infty)\) 递增.

又当 \(x\to 0^+\) 时，\(G(x)\to+\infty\)；当 \(x\to+\infty\) 时，\(G(x)\to+\infty\).
所以 \(f(x)\) 在 \((0,+\infty)\) 只有一个零点时，
\(a=G(2)=\frac{e^2}{4}\).
【法二】\(1^{\odot}\) 当 \(a\le 0\) 时，
\(f(x)=e^x-ax^2>0\)，
所以 \(f(x)\) 在 \((0,+\infty)\) 没有零点.

\(2^{\odot}\) 当 \(a>0\) 时，设函数
\(h(x)=1-ax^2e^{-x}\).
则 \(f(x)\) 在 \((0,+\infty)\) 只有一个零点，等价于 \(h(x)\) 在 \((0,+\infty)\) 只有一个零点.
因为
\(h'(x)=x(x-2)e^{-x}\)，
所以当 \(x\in(0,2)\) 时，\(h'(x)<0\)；当 \(x\in(2,+\infty)\) 时，\(h'(x)>0\).
因此 \(h(x)\) 在 \((0,2)\) 递减，在 \((2,+\infty)\) 递增，于是
\(h(x)_{\min}=h(2)=1-\frac{4a}{e^2}\).

当
\(h(2)<0\)，
即
\(a>\frac{e^2}{4}\)
时，由于 \(h(0)=1\)，且当 \(x>0\) 时，\(e^x>x^2\)，可得
\examdisplaymath{
h(4a)=1-\frac{16a^3}{e^{4a}}=1-\frac{16a^3}{(e^{2a})^2}>1-\frac{16a^3}{(2a)^4}=1-\frac1a>0.
}
所以 \(h(x)\) 在 \((0,+\infty)\) 有 \(2\) 个零点.

当
\(h(2)>0\)，
即
\(a<\frac{e^2}{4}\)
时，\(h(x)\) 在 \((0,+\infty)\) 没有零点.

当
\(h(2)=0\)，
即
\(a=\frac{e^2}{4}\)
时，\(h(x)\) 在 \((0,+\infty)\) 只有一个零点.

综上，
\(a=\frac{e^2}{4}\).
\end{solution}

\begin{problem}
在直角坐标系 \(xOy\) 中，曲线 \(C\) 的参数方程为 \((x=2\cos\theta,\ y=4\sin\theta)\)（\(\theta\)为参数），直线 \(l\) 的参数方程为 \((x=1+t\cos\alpha,\ y=2+t\sin\alpha)\)（\(t\)为参数）.

\begin{enumerate}
\item 求 \(C\) 和 \(l\) 的直角坐标方程；
\item 若曲线 \(C\) 截直线 \(l\) 所得线段的中点坐标为 \((1,2)\)，求 \(l\) 的斜率.
\end{enumerate}

\end{problem}

\begin{solution}
（1）曲线 \(C\) 的参数方程为
\examdisplaycases{}{
x=2\cos\theta,\\
y=4\sin\theta
}
\(\qquad (\theta\text{ 为参数})\)，
转化为直角坐标方程为
\(\frac{y^2}{16}+\frac{x^2}{4}=1\).
直线 \(l\) 的参数方程为
\examdisplaycases{}{
x=1+t\cos\alpha,\\
y=2+t\sin\alpha
}
\qquad (t\text{ 为参数})，
转化为直角坐标方程为
\(x\sin\alpha-y\cos\alpha+2\cos\alpha-\sin\alpha=0\).

（2）把直线 \(l\) 的参数方程代入椭圆的方程，得到
\(\frac{(2+t\sin\alpha)^2}{16}+\frac{(1+t\cos\alpha)^2}{4}=1\).
整理得
\((4\cos^2\alpha+\sin^2\alpha)t^2+(8\cos\alpha+4\sin\alpha)t-8=0\).
设对应参数为 \(t_1\)，\(t_2\)，则
\(t_1+t_2=-\frac{8\cos\alpha+4\sin\alpha}{4\cos^2\alpha+\sin^2\alpha}\).
因为 \((1,2)\) 为中点坐标，

\(1^{\odot}\) 当直线的斜率不存在时，\(x=1\). 此时无解，故舍去.

\(2^{\odot}\) 当直线的斜率存在时，因为 \(t_1\) 和 \(t_2\) 为 \(A\)，\(B\) 对应的参数，所以利用中点坐标公式可知
\(\frac{t_1+t_2}{2}=0\).
因此
\(8\cos\alpha+4\sin\alpha=0\).
解得
\(\tan\alpha=-2\).
即直线 \(l\) 的斜率为 \(-2\).
\end{solution}

\begin{problem}
设函数 \(f(x)=5-|x+a|-|x-2|\).

\begin{enumerate}
\item 当 \(a=1\) 时，求不等式 \(f(x)\ge 0\) 的解集；
\item 若 \(f(x)\le 1\)，求 \(a\) 的取值范围.
\end{enumerate}

\end{problem}

\begin{solution}
（1）当 \(a=1\) 时，
\examdisplaycases{f(x)=5-|x+1|-|x-2|=}{
2x+4, & x\le -1,\\
2, & -1<x<2,\\
-2x+6, & x\ge 2
}

当 \(x\le -1\) 时，
\(f(x)=2x+4\ge 0\)，
解得
\(-2\le x\le -1\).

当 \(-1<x<2\) 时，
\(f(x)=2\ge 0\)
恒成立，即
\(-1<x<2\).

当 \(x\ge 2\) 时，
\(f(x)=-2x+6\ge 0\)，
解得
\(2\le x\le 3\).

综上所述，不等式 \(f(x)\ge 0\) 的解集为
\([-2,3]\).

（2）因为
\(f(x)\le 1\)，
所以
\(5-|x+a|-|x-2|\le 1\)，
即
\(|x+a|+|x-2|\ge 4\).
又
\(|x+a|+|x-2|=|x+a|+|2-x|\ge |x+a+2-x|=|a+2|\).
因此
\(|a+2|\ge 4\).
解得
\(a\le -6\)
或
\(a\ge 2\).
故 \(a\) 的取值范围为
\((-\infty,-6]\cup[2,+\infty)\).
\end{solution}
